\section{奖励模型的评估与部署}

\subsection{Benchmark Stack}

课程强调 reward learning 必须伴随评估体系，不然就像在黑箱里调参。一个有效的 benchmark stack 至少包含三层：

\begin{itemize}
    \item \textbf{自动指标}：log-likelihood、reward model score、KL divergence，与原始 policy 比较是否有所提升
    \item \textbf{人类深度评估}：长轮对话、OpenAI Arena、Anthropic ANAI benchmark 等，通过人工偏好标签持续校准 reward model
    \item \textbf{安全测试}：通过 adversarial probing、red team prompts、low-resource attack attempts，验证 reward model 在 edge case 中仍保持一致
\end{itemize}

\begin{importantbox}{评估的“闭环”原则}
数据 → 奖励 → 策略 → 评估 → 数据。每一轮都要记录日志、复现指标、对比 baseline，并把失败案例用于下一轮 reward model 的训练。否则看似走高的 reward score 很容易是 reward hacking。

\end{importantbox}

\begin{figure}[H]
\centering
\includegraphics[width=0.92\textwidth]{08_cs224r_reward_learning_2025.pdf}
\caption{Reward learning pipeline 的评估反馈回路（来源：lecture slide）}
\end{figure}

\subsection{本章小结}

评估体系必须同时关注自动指标、人类偏好、以及安全触发器，三者结合才能确保 reward learning 模型真正提升体验而非“作弊”。

